%% assessment-template.tex -- starting point for a new randomized assessment.
%%
%% Copy this file, rename it, then edit the three marked spots:
%%   1. the metadata below (\class, \examnum)
%%   2. the question text and figure
%%   3. make_problem() -- the randomized values and the worked solution
%%
%% Build (see README.md):
%%     pdflatex -interaction=nonstopmode assessment-template.tex
%%     pythontex assessment-template.tex
%%     pdflatex -interaction=nonstopmode assessment-template.tex
%%
%% The whole shared preamble lives in jjc-assessment.sty, so this file stays
%% short. Existing assessments predate that package and keep their own inline
%% preambles -- they are unaffected by anything here.

\documentclass[12pt]{exam}
\usepackage{jjc-assessment}

%% ---------------------------------------------------------------- 1. metadata
\renewcommand{\class}{PHYS 202}
\renewcommand{\examnum}{CO00 (v1)}
\renewcommand{\term}{Fall 2026}

\begin{document}

\assessmentheader

\begin{questions}

%% ------------------------------------------------------- 2. question + figure
\question For this circuit:

\begin{figure}[!h]
\begin{center}\begin{circuitikz}\draw
    (0,0) to[battery, invert, l=$V_b$] (0,3)
    to[resistor=$R_1$] (3,3)
    to[resistor=$R_2$] (3,0)
    to[short] (0,0)
;\end{circuitikz}\end{center}
\end{figure}

%% ------------------------------------------------ 3. randomization + solution
\begin{pycode}
from pylab import *
import random

from randassign import RandAssign
ra = RandAssign()


def make_problem():
    # Pick component values from realistic preferred-value series.
    def rand_battery():
        return random.choice([1.5, 3, 4.5, 6, 9, 12, 15, 18])

    def rand_resistor():
        return random.choice([220, 270, 330, 390, 470, 560, 680, 820, 1000])

    R1 = rand_resistor()
    R2 = rand_resistor()
    Vb = rand_battery()

    # Worked solution: series pair, so one current throughout.
    Req = R1 + R2
    I = Vb / Req
    V1 = I * R1
    V2 = I * R2

    # NOTE: doubled braces are .format() escapes for literal LaTeX braces.
    # Use a raw string (r''') so backslashes stay literal.
    problem_template = r'''
    \textbf{{Given}}:
    $V_b=\SI{{{0:.3g}}}{{\volt}}$,
    $R_1=\SI{{{1:.4g}}}{{\ohm}}$, and
    $R_2=\SI{{{2:.4g}}}{{\ohm}}$,
    determine the current through the circuit and the potential difference
    across each resistor.
    '''
    problem = problem_template.format(Vb, R1, R2)

    return problem, I, V1, V2


problem, I, V1, V2 = make_problem()

# Record the answer key. randassign collects these into randassign/solutions/.
solution = (r'$I=\SI{{{0:.3g}}}{{\ampere}}$, '
            r'$V_1=\SI{{{1:.3g}}}{{\volt}}$, '
            r'$V_2=\SI{{{2:.3g}}}{{\volt}}$').format(I, V1, V2)
ra.addsoln(solution)

print(problem)
\end{pycode}

\vspace{2in}

\end{questions}

\end{document}
